\chapter{Intrauterine Insemination}

\section*{Overview}
Intrauterine insemination (IUI) is a fertility treatment that places concentrated, washed sperm directly into the uterine cavity around the time of ovulation. It brings sperm closer to the egg to increase the chance of fertilization.

\section*{Alternative Names}
IUI, artificial insemination, donor insemination, sperm insemination

\section*{Principle}
A semen sample is washed and concentrated to remove seminal fluid and select motile sperm. The prepared sample is injected through a thin catheter across the cervix into the uterus near ovulation. This shortens the distance sperm must travel and overcomes mild cervical or semen factors impairing natural conception.

\section*{History}
Early attempts at artificial insemination date to the 18th century. Modern intrauterine insemination with washed sperm became established in the 1980s and remains a first-line fertility treatment.

\section*{Why It Is Done}
It is ordered for unexplained infertility, mild male factor infertility, cervical mucus abnormalities, or ejaculatory dysfunction. It is also used with donor sperm for single women and same-sex female couples, and may be combined with ovulation induction in women with anovulation.

\section*{Is This a Primary or Secondary Test or Procedure}
It is a primary, low-complexity treatment for mild fertility problems. It may be secondary when used after ovulation induction or as an alternative to more invasive treatments such as in vitro fertilization.

\section*{How It Is Performed}
Ovulation timing is tracked with ultrasound and sometimes urine luteinizing hormone tests, with or without ovulation induction medication. The male partner provides a semen sample by masturbation, which is washed and concentrated in the laboratory. The prepared sperm is injected through a thin, soft catheter into the uterus; the procedure takes minutes and requires no anesthesia.

\section*{Preparation}
Ovulation is monitored with serial ultrasounds, and the cycle may be coordinated with ovulation induction medication. The partner abstains from ejaculation for 2-5 days before providing the semen sample. A full bladder is sometimes requested before the procedure.

\section*{Contraindications and Precautions}
Contraindications include tubal obstruction, active pelvic infection, and severe male factor infertility. Precautions include screening for sexually transmitted infections in both partners and counseling about the risk of multiple gestation when ovulation induction is used.

\section*{Risks}
Risks are minimal and include mild cramping, spotting, and a small risk of infection. Ovarian hyperstimulation syndrome and multiple pregnancy may occur when ovulation induction medication is combined with IUI.

\section*{Aftercare and Recovery}
The patient rests for 10-15 minutes and may resume normal activities, avoiding strenuous exercise for the day. A pregnancy test is performed about 2 weeks later.

\section*{Outcome and Follow-up}
A positive pregnancy test approximately 2 weeks after insemination indicates conception. If the cycle is unsuccessful, the fertility specialist reviews the response and may adjust medications or consider additional cycles or more advanced treatment.

\section*{Expected Results}
A properly timed insemination with an adequate number of motile sperm delivered into the uterine cavity and a subsequent positive pregnancy test.

\section*{Follow-up}
A serum hCG pregnancy test is done about 2 weeks after the procedure, and an early ultrasound confirms viability if positive. Unsuccessful cycles may be followed by repeat IUI or progression to in vitro fertilization.

\section*{When to Seek Medical Care}
Follow the procedure team's specific instructions. Seek urgent medical assessment for severe or worsening pain, uncontrolled bleeding, fever or signs of infection, trouble breathing, chest pain, fainting, new weakness or confusion, inability to urinate, or any rapidly worsening symptom. The appropriate warning signs depend on the procedure and the patient's medical condition.

\section*{References}
\begin{enumerate}
  \item MedlinePlus. Intrauterine Insemination (IUI). U.S. National Library of Medicine; 2025.
  \item Current Medical Diagnosis and Treatment. McGraw-Hill; 2025.
\end{enumerate}

\clearpage
